% paper72-agent-verification-loop.tex — Verified Agent Actions via Sheaf-Theoretic Certificates
%   Paper 72 of the Judgment Geometry series.
% Compile: pdflatex paper72-agent-verification-loop
\documentclass[11pt]{article}
\usepackage{lmodern}
% jugeo-common.tex — shared preamble for all Judgment Geometry papers
% Include via: % jugeo-common.tex — shared preamble for all Judgment Geometry papers
% Include via: % jugeo-common.tex — shared preamble for all Judgment Geometry papers
% Include via: \input{jugeo-common}

% ── Packages ────────────────────────────────────────────────────────────
\usepackage{amsmath,amssymb,amsthm,mathtools}
\usepackage{tikz-cd}
\usepackage{listings}
\usepackage{xcolor}
\usepackage{xspace}
\usepackage{booktabs}
\usepackage{multirow}
\usepackage{enumitem}
\usepackage[expansion=false]{microtype}
\usepackage{hyperref}
\usepackage{cleveref}
% \usepackage{thmtools}  % disabled: not available in this TeX installation
\usepackage{float}
\usepackage{caption}
\usepackage{subcaption}
\usepackage{tabularx}
\usepackage{stmaryrd}       % \llbracket, \rrbracket
\usepackage{bbm}            % \mathbbm{1}

% ── Page geometry ───────────────────────────────────────────────────────
\usepackage[margin=1in]{geometry}

% ── Theorem environments ────────────────────────────────────────────────
\theoremstyle{plain}
\newtheorem{theorem}{Theorem}[section]
\newtheorem{lemma}[theorem]{Lemma}
\newtheorem{proposition}[theorem]{Proposition}
\newtheorem{corollary}[theorem]{Corollary}

\theoremstyle{definition}
\newtheorem{definition}[theorem]{Definition}
\newtheorem{example}[theorem]{Example}
\newtheorem{construction}[theorem]{Construction}

\theoremstyle{remark}
\newtheorem{remark}[theorem]{Remark}
\newtheorem{notation}[theorem]{Notation}

% ── Python listings style ──────────────────────────────────────────────
\definecolor{jugeoBlue}{HTML}{2563EB}
\definecolor{jugeoGreen}{HTML}{059669}
\definecolor{jugeoRed}{HTML}{DC2626}
\definecolor{jugeoPurple}{HTML}{7C3AED}
\definecolor{jugeoGray}{HTML}{6B7280}
\definecolor{jugeoBg}{HTML}{F8FAFC}

\lstdefinestyle{jugeo-python}{
  language=Python,
  basicstyle=\ttfamily\small,
  keywordstyle=\color{jugeoBlue}\bfseries,
  stringstyle=\color{jugeoGreen},
  commentstyle=\color{jugeoGray}\itshape,
  numberstyle=\tiny\color{jugeoGray},
  backgroundcolor=\color{jugeoBg},
  numbers=left,
  numbersep=6pt,
  frame=single,
  framerule=0.4pt,
  rulecolor=\color{jugeoGray!40},
  breaklines=true,
  breakatwhitespace=true,
  showstringspaces=false,
  tabsize=4,
  xleftmargin=1.5em,
  framexleftmargin=1.5em,
  morekeywords={dataclass,frozen,field,Enum,IntEnum,auto,Any,
                Mapping,Sequence,Optional,match,case},
  literate={→}{{$\to$}}1 {←}{{$\leftarrow$}}1
           {≤}{{$\leq$}}1 {≥}{{$\geq$}}1 {≠}{{$\neq$}}1
           {∈}{{$\in$}}1 {∀}{{$\forall$}}1 {∃}{{$\exists$}}1
           {⊕}{{$\oplus$}}1 {⊖}{{$\ominus$}}1,
  escapeinside={(*@}{@*)},
}
\lstset{style=jugeo-python}

\lstdefinestyle{jugeo-output}{
  basicstyle=\ttfamily\small,
  backgroundcolor=\color{jugeoBg},
  frame=single,
  framerule=0.4pt,
  rulecolor=\color{jugeoGray!40},
  breaklines=true,
  showstringspaces=false,
  xleftmargin=1.5em,
  framexleftmargin=0.5em,
  numbers=none,
}

% ── JuGeo-specific macros ──────────────────────────────────────────────

% Judgment 8-tuple
\newcommand{\J}{\mathcal{J}}
\newcommand{\Jud}[8]{(#1,\,#2,\,#3,\,#4,\,#5,\,#6,\,#7,\,#8)}
\newcommand{\JudShort}[1]{\J_{#1}}

% Components
\newcommand{\coord}{\mathsf{c}}            % coordinate
\newcommand{\prop}{\varphi}                 % proposition
\newcommand{\carrier}{\mathcal{A}}          % carrier type
\newcommand{\evid}{\mathcal{E}}             % evidence bundle
\newcommand{\oblig}{\mathcal{O}}            % obligations
\newcommand{\obstr}{\mathcal{B}}            % obstructions
\newcommand{\trust}{\mathcal{T}}            % trust annotation
\newcommand{\prov}{\Pi}                     % provenance

% Trust algebra
\newcommand{\Talg}{\mathfrak{T}}            % trust ordered algebra
\newcommand{\Eadm}{\mathcal{E}_{\mathrm{adm}}}
\newcommand{\tleq}{\preceq}                 % trust ordering
\newcommand{\tjoin}{\oplus}                 % conservative join
\newcommand{\tatten}{\ominus}               % attenuation
\newcommand{\tpromote}{\uparrow_\pi}        % promotion
\newcommand{\tdemote}{\downarrow_\chi}       % demotion

% Trust tiers
\newcommand{\Tcontra}{\ensuremath{\mathsf{CONTRADICTED}}}
\newcommand{\Tunver}{\ensuremath{\mathsf{UNVERIFIED}}}
\newcommand{\Tcopilot}{\ensuremath{\mathsf{COPILOT}}}
\newcommand{\Toracle}{\ensuremath{\mathsf{ORACLE}}}
\newcommand{\Truntime}{\ensuremath{\mathsf{RUNTIME}}}
\newcommand{\Tsolver}{\ensuremath{\mathsf{SOLVER}}}
\newcommand{\Tproof}{\ensuremath{\mathsf{PROOF}}}

% Site and geometry
\newcommand{\Site}{\mathcal{S}}             % semantic site
\newcommand{\Cov}{\mathrm{Cov}}            % covering family
\newcommand{\Ob}{\mathrm{Ob}}              % objects
\newcommand{\Mor}{\mathrm{Mor}}            % morphisms
\newcommand{\res}{\mathrm{res}}            % restriction
\newcommand{\Cech}{\check{C}}              % Čech complex
\newcommand{\Hcoh}[1]{H^{#1}}             % cohomology group

% Descent
\newcommand{\descent}{\mathrm{desc}}
\newcommand{\glue}{\mathrm{glue}}
\newcommand{\overlap}[2]{#1 \cap #2}

% Semantic moves
\newcommand{\VERIFY}{\textsc{Verify}}
\newcommand{\CONSTRUCT}{\textsc{Construct}}
\newcommand{\NORMALIZE}{\textsc{Normalize}}
\newcommand{\GLUE}{\textsc{Glue}}
\newcommand{\RESTRICT}{\textsc{Restrict}}
\newcommand{\TRANSPORT}{\textsc{Transport}}
\newcommand{\REPAIR}{\textsc{Repair}}
\newcommand{\PROMOTE}{\textsc{Promote}}
\newcommand{\CHALLENGE}{\textsc{Challenge}}
\newcommand{\ESCALATE}{\textsc{Escalate}}

% Fragments
\newcommand{\QFLIA}{\texttt{QF\_LIA}}
\newcommand{\QFLRA}{\texttt{QF\_LRA}}
\newcommand{\QFBV}{\texttt{QF\_BV}}
\newcommand{\QFUF}{\texttt{QF\_UF}}

% Misc
\newcommand{\Python}{\textsc{Python}}
\newcommand{\JuGeo}{\textsc{JuGeo}}
\newcommand{\LEAN}{\textsc{Lean}\,4}
\newcommand{\Fstar}{F$^{\star}$}
\newcommand{\Coq}{\textsc{Coq}}
\newcommand{\Dafny}{\textsc{Dafny}}

% Common shortcuts
\newcommand{\ie}{i.e.\@\xspace}
\newcommand{\eg}{e.g.\@\xspace}
\newcommand{\cf}{cf.\@\xspace}
\newcommand{\wrt}{w.r.t.\@\xspace}

% Evaluation shortcuts
\newcommand{\numcases}{300}
\newcommand{\numspec}{100}
\newcommand{\numeq}{100}
\newcommand{\numbug}{100}
\newcommand{\medtime}{6.5\,ms}
\newcommand{\totaltime}{1.949\,s}


% ── Packages ────────────────────────────────────────────────────────────
\usepackage{amsmath,amssymb,amsthm,mathtools}
\usepackage{tikz-cd}
\usepackage{listings}
\usepackage{xcolor}
\usepackage{xspace}
\usepackage{booktabs}
\usepackage{multirow}
\usepackage{enumitem}
\usepackage[expansion=false]{microtype}
\usepackage{hyperref}
\usepackage{cleveref}
% \usepackage{thmtools}  % disabled: not available in this TeX installation
\usepackage{float}
\usepackage{caption}
\usepackage{subcaption}
\usepackage{tabularx}
\usepackage{stmaryrd}       % \llbracket, \rrbracket
\usepackage{bbm}            % \mathbbm{1}

% ── Page geometry ───────────────────────────────────────────────────────
\usepackage[margin=1in]{geometry}

% ── Theorem environments ────────────────────────────────────────────────
\theoremstyle{plain}
\newtheorem{theorem}{Theorem}[section]
\newtheorem{lemma}[theorem]{Lemma}
\newtheorem{proposition}[theorem]{Proposition}
\newtheorem{corollary}[theorem]{Corollary}

\theoremstyle{definition}
\newtheorem{definition}[theorem]{Definition}
\newtheorem{example}[theorem]{Example}
\newtheorem{construction}[theorem]{Construction}

\theoremstyle{remark}
\newtheorem{remark}[theorem]{Remark}
\newtheorem{notation}[theorem]{Notation}

% ── Python listings style ──────────────────────────────────────────────
\definecolor{jugeoBlue}{HTML}{2563EB}
\definecolor{jugeoGreen}{HTML}{059669}
\definecolor{jugeoRed}{HTML}{DC2626}
\definecolor{jugeoPurple}{HTML}{7C3AED}
\definecolor{jugeoGray}{HTML}{6B7280}
\definecolor{jugeoBg}{HTML}{F8FAFC}

\lstdefinestyle{jugeo-python}{
  language=Python,
  basicstyle=\ttfamily\small,
  keywordstyle=\color{jugeoBlue}\bfseries,
  stringstyle=\color{jugeoGreen},
  commentstyle=\color{jugeoGray}\itshape,
  numberstyle=\tiny\color{jugeoGray},
  backgroundcolor=\color{jugeoBg},
  numbers=left,
  numbersep=6pt,
  frame=single,
  framerule=0.4pt,
  rulecolor=\color{jugeoGray!40},
  breaklines=true,
  breakatwhitespace=true,
  showstringspaces=false,
  tabsize=4,
  xleftmargin=1.5em,
  framexleftmargin=1.5em,
  morekeywords={dataclass,frozen,field,Enum,IntEnum,auto,Any,
                Mapping,Sequence,Optional,match,case},
  literate={→}{{$\to$}}1 {←}{{$\leftarrow$}}1
           {≤}{{$\leq$}}1 {≥}{{$\geq$}}1 {≠}{{$\neq$}}1
           {∈}{{$\in$}}1 {∀}{{$\forall$}}1 {∃}{{$\exists$}}1
           {⊕}{{$\oplus$}}1 {⊖}{{$\ominus$}}1,
  escapeinside={(*@}{@*)},
}
\lstset{style=jugeo-python}

\lstdefinestyle{jugeo-output}{
  basicstyle=\ttfamily\small,
  backgroundcolor=\color{jugeoBg},
  frame=single,
  framerule=0.4pt,
  rulecolor=\color{jugeoGray!40},
  breaklines=true,
  showstringspaces=false,
  xleftmargin=1.5em,
  framexleftmargin=0.5em,
  numbers=none,
}

% ── JuGeo-specific macros ──────────────────────────────────────────────

% Judgment 8-tuple
\newcommand{\J}{\mathcal{J}}
\newcommand{\Jud}[8]{(#1,\,#2,\,#3,\,#4,\,#5,\,#6,\,#7,\,#8)}
\newcommand{\JudShort}[1]{\J_{#1}}

% Components
\newcommand{\coord}{\mathsf{c}}            % coordinate
\newcommand{\prop}{\varphi}                 % proposition
\newcommand{\carrier}{\mathcal{A}}          % carrier type
\newcommand{\evid}{\mathcal{E}}             % evidence bundle
\newcommand{\oblig}{\mathcal{O}}            % obligations
\newcommand{\obstr}{\mathcal{B}}            % obstructions
\newcommand{\trust}{\mathcal{T}}            % trust annotation
\newcommand{\prov}{\Pi}                     % provenance

% Trust algebra
\newcommand{\Talg}{\mathfrak{T}}            % trust ordered algebra
\newcommand{\Eadm}{\mathcal{E}_{\mathrm{adm}}}
\newcommand{\tleq}{\preceq}                 % trust ordering
\newcommand{\tjoin}{\oplus}                 % conservative join
\newcommand{\tatten}{\ominus}               % attenuation
\newcommand{\tpromote}{\uparrow_\pi}        % promotion
\newcommand{\tdemote}{\downarrow_\chi}       % demotion

% Trust tiers
\newcommand{\Tcontra}{\ensuremath{\mathsf{CONTRADICTED}}}
\newcommand{\Tunver}{\ensuremath{\mathsf{UNVERIFIED}}}
\newcommand{\Tcopilot}{\ensuremath{\mathsf{COPILOT}}}
\newcommand{\Toracle}{\ensuremath{\mathsf{ORACLE}}}
\newcommand{\Truntime}{\ensuremath{\mathsf{RUNTIME}}}
\newcommand{\Tsolver}{\ensuremath{\mathsf{SOLVER}}}
\newcommand{\Tproof}{\ensuremath{\mathsf{PROOF}}}

% Site and geometry
\newcommand{\Site}{\mathcal{S}}             % semantic site
\newcommand{\Cov}{\mathrm{Cov}}            % covering family
\newcommand{\Ob}{\mathrm{Ob}}              % objects
\newcommand{\Mor}{\mathrm{Mor}}            % morphisms
\newcommand{\res}{\mathrm{res}}            % restriction
\newcommand{\Cech}{\check{C}}              % Čech complex
\newcommand{\Hcoh}[1]{H^{#1}}             % cohomology group

% Descent
\newcommand{\descent}{\mathrm{desc}}
\newcommand{\glue}{\mathrm{glue}}
\newcommand{\overlap}[2]{#1 \cap #2}

% Semantic moves
\newcommand{\VERIFY}{\textsc{Verify}}
\newcommand{\CONSTRUCT}{\textsc{Construct}}
\newcommand{\NORMALIZE}{\textsc{Normalize}}
\newcommand{\GLUE}{\textsc{Glue}}
\newcommand{\RESTRICT}{\textsc{Restrict}}
\newcommand{\TRANSPORT}{\textsc{Transport}}
\newcommand{\REPAIR}{\textsc{Repair}}
\newcommand{\PROMOTE}{\textsc{Promote}}
\newcommand{\CHALLENGE}{\textsc{Challenge}}
\newcommand{\ESCALATE}{\textsc{Escalate}}

% Fragments
\newcommand{\QFLIA}{\texttt{QF\_LIA}}
\newcommand{\QFLRA}{\texttt{QF\_LRA}}
\newcommand{\QFBV}{\texttt{QF\_BV}}
\newcommand{\QFUF}{\texttt{QF\_UF}}

% Misc
\newcommand{\Python}{\textsc{Python}}
\newcommand{\JuGeo}{\textsc{JuGeo}}
\newcommand{\LEAN}{\textsc{Lean}\,4}
\newcommand{\Fstar}{F$^{\star}$}
\newcommand{\Coq}{\textsc{Coq}}
\newcommand{\Dafny}{\textsc{Dafny}}

% Common shortcuts
\newcommand{\ie}{i.e.\@\xspace}
\newcommand{\eg}{e.g.\@\xspace}
\newcommand{\cf}{cf.\@\xspace}
\newcommand{\wrt}{w.r.t.\@\xspace}

% Evaluation shortcuts
\newcommand{\numcases}{300}
\newcommand{\numspec}{100}
\newcommand{\numeq}{100}
\newcommand{\numbug}{100}
\newcommand{\medtime}{6.5\,ms}
\newcommand{\totaltime}{1.949\,s}


% ── Packages ────────────────────────────────────────────────────────────
\usepackage{amsmath,amssymb,amsthm,mathtools}
\usepackage{tikz-cd}
\usepackage{listings}
\usepackage{xcolor}
\usepackage{xspace}
\usepackage{booktabs}
\usepackage{multirow}
\usepackage{enumitem}
\usepackage[expansion=false]{microtype}
\usepackage{hyperref}
\usepackage{cleveref}
% \usepackage{thmtools}  % disabled: not available in this TeX installation
\usepackage{float}
\usepackage{caption}
\usepackage{subcaption}
\usepackage{tabularx}
\usepackage{stmaryrd}       % \llbracket, \rrbracket
\usepackage{bbm}            % \mathbbm{1}

% ── Page geometry ───────────────────────────────────────────────────────
\usepackage[margin=1in]{geometry}

% ── Theorem environments ────────────────────────────────────────────────
\theoremstyle{plain}
\newtheorem{theorem}{Theorem}[section]
\newtheorem{lemma}[theorem]{Lemma}
\newtheorem{proposition}[theorem]{Proposition}
\newtheorem{corollary}[theorem]{Corollary}

\theoremstyle{definition}
\newtheorem{definition}[theorem]{Definition}
\newtheorem{example}[theorem]{Example}
\newtheorem{construction}[theorem]{Construction}

\theoremstyle{remark}
\newtheorem{remark}[theorem]{Remark}
\newtheorem{notation}[theorem]{Notation}

% ── Python listings style ──────────────────────────────────────────────
\definecolor{jugeoBlue}{HTML}{2563EB}
\definecolor{jugeoGreen}{HTML}{059669}
\definecolor{jugeoRed}{HTML}{DC2626}
\definecolor{jugeoPurple}{HTML}{7C3AED}
\definecolor{jugeoGray}{HTML}{6B7280}
\definecolor{jugeoBg}{HTML}{F8FAFC}

\lstdefinestyle{jugeo-python}{
  language=Python,
  basicstyle=\ttfamily\small,
  keywordstyle=\color{jugeoBlue}\bfseries,
  stringstyle=\color{jugeoGreen},
  commentstyle=\color{jugeoGray}\itshape,
  numberstyle=\tiny\color{jugeoGray},
  backgroundcolor=\color{jugeoBg},
  numbers=left,
  numbersep=6pt,
  frame=single,
  framerule=0.4pt,
  rulecolor=\color{jugeoGray!40},
  breaklines=true,
  breakatwhitespace=true,
  showstringspaces=false,
  tabsize=4,
  xleftmargin=1.5em,
  framexleftmargin=1.5em,
  morekeywords={dataclass,frozen,field,Enum,IntEnum,auto,Any,
                Mapping,Sequence,Optional,match,case},
  literate={→}{{$\to$}}1 {←}{{$\leftarrow$}}1
           {≤}{{$\leq$}}1 {≥}{{$\geq$}}1 {≠}{{$\neq$}}1
           {∈}{{$\in$}}1 {∀}{{$\forall$}}1 {∃}{{$\exists$}}1
           {⊕}{{$\oplus$}}1 {⊖}{{$\ominus$}}1,
  escapeinside={(*@}{@*)},
}
\lstset{style=jugeo-python}

\lstdefinestyle{jugeo-output}{
  basicstyle=\ttfamily\small,
  backgroundcolor=\color{jugeoBg},
  frame=single,
  framerule=0.4pt,
  rulecolor=\color{jugeoGray!40},
  breaklines=true,
  showstringspaces=false,
  xleftmargin=1.5em,
  framexleftmargin=0.5em,
  numbers=none,
}

% ── JuGeo-specific macros ──────────────────────────────────────────────

% Judgment 8-tuple
\newcommand{\J}{\mathcal{J}}
\newcommand{\Jud}[8]{(#1,\,#2,\,#3,\,#4,\,#5,\,#6,\,#7,\,#8)}
\newcommand{\JudShort}[1]{\J_{#1}}

% Components
\newcommand{\coord}{\mathsf{c}}            % coordinate
\newcommand{\prop}{\varphi}                 % proposition
\newcommand{\carrier}{\mathcal{A}}          % carrier type
\newcommand{\evid}{\mathcal{E}}             % evidence bundle
\newcommand{\oblig}{\mathcal{O}}            % obligations
\newcommand{\obstr}{\mathcal{B}}            % obstructions
\newcommand{\trust}{\mathcal{T}}            % trust annotation
\newcommand{\prov}{\Pi}                     % provenance

% Trust algebra
\newcommand{\Talg}{\mathfrak{T}}            % trust ordered algebra
\newcommand{\Eadm}{\mathcal{E}_{\mathrm{adm}}}
\newcommand{\tleq}{\preceq}                 % trust ordering
\newcommand{\tjoin}{\oplus}                 % conservative join
\newcommand{\tatten}{\ominus}               % attenuation
\newcommand{\tpromote}{\uparrow_\pi}        % promotion
\newcommand{\tdemote}{\downarrow_\chi}       % demotion

% Trust tiers
\newcommand{\Tcontra}{\ensuremath{\mathsf{CONTRADICTED}}}
\newcommand{\Tunver}{\ensuremath{\mathsf{UNVERIFIED}}}
\newcommand{\Tcopilot}{\ensuremath{\mathsf{COPILOT}}}
\newcommand{\Toracle}{\ensuremath{\mathsf{ORACLE}}}
\newcommand{\Truntime}{\ensuremath{\mathsf{RUNTIME}}}
\newcommand{\Tsolver}{\ensuremath{\mathsf{SOLVER}}}
\newcommand{\Tproof}{\ensuremath{\mathsf{PROOF}}}

% Site and geometry
\newcommand{\Site}{\mathcal{S}}             % semantic site
\newcommand{\Cov}{\mathrm{Cov}}            % covering family
\newcommand{\Ob}{\mathrm{Ob}}              % objects
\newcommand{\Mor}{\mathrm{Mor}}            % morphisms
\newcommand{\res}{\mathrm{res}}            % restriction
\newcommand{\Cech}{\check{C}}              % Čech complex
\newcommand{\Hcoh}[1]{H^{#1}}             % cohomology group

% Descent
\newcommand{\descent}{\mathrm{desc}}
\newcommand{\glue}{\mathrm{glue}}
\newcommand{\overlap}[2]{#1 \cap #2}

% Semantic moves
\newcommand{\VERIFY}{\textsc{Verify}}
\newcommand{\CONSTRUCT}{\textsc{Construct}}
\newcommand{\NORMALIZE}{\textsc{Normalize}}
\newcommand{\GLUE}{\textsc{Glue}}
\newcommand{\RESTRICT}{\textsc{Restrict}}
\newcommand{\TRANSPORT}{\textsc{Transport}}
\newcommand{\REPAIR}{\textsc{Repair}}
\newcommand{\PROMOTE}{\textsc{Promote}}
\newcommand{\CHALLENGE}{\textsc{Challenge}}
\newcommand{\ESCALATE}{\textsc{Escalate}}

% Fragments
\newcommand{\QFLIA}{\texttt{QF\_LIA}}
\newcommand{\QFLRA}{\texttt{QF\_LRA}}
\newcommand{\QFBV}{\texttt{QF\_BV}}
\newcommand{\QFUF}{\texttt{QF\_UF}}

% Misc
\newcommand{\Python}{\textsc{Python}}
\newcommand{\JuGeo}{\textsc{JuGeo}}
\newcommand{\LEAN}{\textsc{Lean}\,4}
\newcommand{\Fstar}{F$^{\star}$}
\newcommand{\Coq}{\textsc{Coq}}
\newcommand{\Dafny}{\textsc{Dafny}}

% Common shortcuts
\newcommand{\ie}{i.e.\@\xspace}
\newcommand{\eg}{e.g.\@\xspace}
\newcommand{\cf}{cf.\@\xspace}
\newcommand{\wrt}{w.r.t.\@\xspace}

% Evaluation shortcuts
\newcommand{\numcases}{300}
\newcommand{\numspec}{100}
\newcommand{\numeq}{100}
\newcommand{\numbug}{100}
\newcommand{\medtime}{6.5\,ms}
\newcommand{\totaltime}{1.949\,s}

% \input{data-paper72}  % data file does not exist: no measured data for paper 72

% ── Additional packages ────────────────────────────────────────────
\usepackage{array}
\usepackage{tikz}
\usetikzlibrary{arrows.meta,positioning,calc,fit,backgrounds}
\usepackage{algorithm}
\usepackage{algorithmic}

% ── Paper-specific macros ──────────────────────────────────────────
\newcommand{\ActionSite}{\mathcal{A}_{\mathrm{site}}}
\newcommand{\ActionSheaf}{\mathscr{A}}
\newcommand{\CertBundle}{\mathcal{K}}
\newcommand{\AgentOrch}{\textsc{AgentOrchestrator}}
\newcommand{\CertGen}{\textsc{CertificateGenerator}}
\newcommand{\ActionVerif}{\textsc{ActionVerifier}}
\newcommand{\GateKeeper}{\textsc{GateKeeper}}
\newcommand{\ActionSpace}{\mathcal{A}}
\newcommand{\ActionType}{\tau_a}
\newcommand{\PreCond}{\mathrm{pre}}
\newcommand{\PostCond}{\mathrm{post}}
\newcommand{\Effect}{\mathrm{eff}}
\newcommand{\SafeAction}{\mathsf{SAFE}}
\newcommand{\UnsafeAction}{\mathsf{UNSAFE}}
\newcommand{\PendingAction}{\mathsf{PENDING}}
\newcommand{\CertValid}{\mathsf{VALID}}
\newcommand{\CertRevoked}{\mathsf{REVOKED}}

\title{\textbf{Verified Agent Actions via Sheaf-Theoretic Certificates}}

\author{JuGeo Research Group}
\date{\today}

\begin{document}
\maketitle

% ─────────────────────────────────────────────────────────────────────
\begin{abstract}
We present a framework for formally verifying the actions of AI code
agents before execution, preventing hallucinated or incorrect code from
reaching production.  The agent's action space is modeled as a
\emph{semantic site} $\ActionSite$ where objects are action types (file
writes, function calls, API requests, shell commands, database queries),
morphisms capture dependencies between actions, and covering families
encode action decomposition.  Each action produces a \emph{local section}
of the action sheaf $\ActionSheaf$, carrying preconditions, postconditions,
and effects.  Before committing any action, the $\GateKeeper$ orchestrator
runs \JuGeo's descent engine to verify that the action's section is
compatible with previously committed actions---i.e., that the local
sections \emph{glue} into a consistent global state.  When gluing fails,
the obstruction class identifies the specific conflict (type mismatch,
state violation, API incompatibility), enabling targeted repair.
We describe an evaluation methodology for measuring hallucination catch
rates, bug reduction, and verification overhead, with quantitative
validation left to future work.
\end{abstract}

\newpage

% =====================================================================
\section{Introduction}\label{sec:intro}
% =====================================================================

AI coding agents---systems that autonomously write, edit, and execute
code---are increasingly deployed in software development workflows.
GitHub Copilot, Claude Code, Devin, and similar agents generate code
from natural language instructions, call APIs, execute shell commands,
and modify files.  However, these agents share a fundamental
limitation: they can hallucinate---generating code that is
syntactically plausible but semantically incorrect, unsafe, or
inconsistent with the existing codebase.

The consequences of unverified agent actions range from subtle bugs
to catastrophic failures.  An agent might generate a function that
violates an implicit invariant maintained by other modules, call an
API with incorrect parameters, or write to a file in a way that
conflicts with concurrent operations.  Current mitigation strategies
(code review, linting, testing) are reactive: they detect problems
\emph{after} the agent has acted.

We propose a \emph{proactive} approach: verify each agent action
\emph{before} execution using \JuGeo's sheaf-theoretic verification
framework.  The key insight is that the agent's action space naturally
forms a \emph{site}---a topological structure where actions are
objects, dependencies are morphisms, and covering families capture
how complex actions decompose into simpler ones.  Each action produces
a \emph{local section} of the action sheaf, carrying formal metadata:
preconditions, postconditions, effects, and type constraints.  The
descent engine checks whether a new action's section is compatible
with the existing committed sections---whether they \emph{glue} into
a consistent global state.

\paragraph{Contributions.}
\begin{itemize}
  \item A formal model of agent action spaces as semantic sites, with
    a sheaf capturing action semantics (\S\ref{sec:action-site}).
  \item A verify-before-commit protocol using descent
    (\S\ref{sec:verification-protocol}).
  \item Certificate-based trust tracking with revocation
    (\S\ref{sec:certificates}).
  \item Obstruction-based conflict diagnosis with automatic repair
    (\S\ref{sec:obstruction-repair}).
  \item An evaluation methodology for measuring bug
    reduction (\S\ref{sec:experiments}).
\end{itemize}


% =====================================================================
\section{Action Spaces as Semantic Sites}\label{sec:action-site}
% =====================================================================

We formalize the agent's action space as a semantic site, enabling
sheaf-theoretic verification.

% ─────────────────────────────────────────────────────────────────────
\subsection{The Action Site}\label{subsec:action-site-def}

\begin{definition}[Action Type]\label{def:action-type}
An \emph{action type} is a triple
$\ActionType = (\PreCond, \PostCond, \Effect)$ where:
\begin{itemize}
  \item $\PreCond$ is a set of preconditions (assertions about the state
    before execution);
  \item $\PostCond$ is a set of postconditions (assertions about the state
    after execution);
  \item $\Effect$ is a set of effects (state modifications produced by
    the action).
\end{itemize}
The canonical action types are: $\mathsf{FileWrite}$, $\mathsf{FuncCall}$,
$\mathsf{APICall}$, $\mathsf{ShellExec}$, and $\mathsf{DBQuery}$.
\end{definition}

\begin{definition}[Action Site]\label{def:action-site-formal}
The \emph{action site} $\ActionSite = (\mathcal{C}_A, \Mor_A, J_A)$ where:
\begin{enumerate}[label=(\alph*)]
  \item $\mathcal{C}_A$ is the category whose objects are action instances
    $a = (\ActionType, \mathrm{params}, \mathrm{target})$;
  \item $\Mor_A(a, b)$ consists of \emph{dependency morphisms}: $a \to b$
    if action $b$ depends on the output of action $a$ (data flow,
    control flow, or resource dependency);
  \item $J_A$ is the Grothendieck topology where a family
    $\{a_i \to a\}$ covers $a$ if executing all $a_i$ in the correct
    order is equivalent to executing $a$.
\end{enumerate}
\end{definition}

\begin{definition}[Action Sheaf]\label{def:action-sheaf}
The \emph{action sheaf} $\ActionSheaf$ on $\ActionSite$ assigns to each
action $a$ the set $\ActionSheaf(a)$ of \emph{valid execution states}---
pairs $(\PreCond_a, \PostCond_a)$ such that executing $a$ in a state
satisfying $\PreCond_a$ yields a state satisfying $\PostCond_a$.
Restriction maps encode how validity propagates along dependencies:
if $f : a \to b$, then $\res_f : \ActionSheaf(b) \to \ActionSheaf(a)$
restricts the validity of $b$ to the preconditions it imposes on $a$.
\end{definition}

\begin{example}[File Write and Function Call]\label{ex:file-func}
Consider an agent that writes a function definition to a file (action
$a_1 : \mathsf{FileWrite}$) and then calls that function (action
$a_2 : \mathsf{FuncCall}$).  The dependency $a_1 \to a_2$ captures
that $a_2$ depends on $a_1$.  The restriction $\res_{a_1 \to a_2}$
verifies the written function's signature matches the call site.
\end{example}

% ─────────────────────────────────────────────────────────────────────
\subsection{Local Sections for Agent Actions}\label{subsec:local-sections}

\begin{definition}[Action Section]\label{def:action-section}
A \emph{local section} for action $a$ is an element
$s_a \in \ActionSheaf(a)$ consisting of:
\begin{enumerate}[label=(\roman*)]
  \item The code or command to be executed;
  \item Formal preconditions derived from the current state;
  \item Expected postconditions and effects;
  \item Type annotations and API contract assertions;
  \item A trust level $t_a \in \Talg$ (initially $\Tcopilot$ for
    LLM-generated actions).
\end{enumerate}
\end{definition}

\begin{theorem}[Section Composition]\label{thm:section-compose}
Let $a_1 \to a_2$ be a dependency morphism with sections
$s_1 \in \ActionSheaf(a_1)$ and $s_2 \in \ActionSheaf(a_2)$.
The \emph{composed section} $s_{1;2}$ satisfies:
\begin{enumerate}[label=(\roman*)]
  \item $\PreCond(s_{1;2}) = \PreCond(s_1)$;
  \item $\PostCond(s_{1;2}) = \PostCond(s_2)$;
  \item $\Effect(s_{1;2}) = \Effect(s_1) \cup \Effect(s_2)$;
  \item $\trust(s_{1;2}) = \trust(s_1) \sqcap \trust(s_2)$.
\end{enumerate}
provided $\PostCond(s_1) \Rightarrow \PreCond(s_2)$ (compatibility).
\end{theorem}

\begin{proof}
The composed section executes $a_1$ then $a_2$.  Condition~(i) is
immediate: the initial precondition is that of $a_1$.  For~(ii), the
final postcondition is that of $a_2$, since $a_2$ executes last.
For~(iii), effects accumulate.  For~(iv), trust is the meet of component
trusts by the trust algebra axioms~\cite{JuGeo2024}.  The compatibility
condition $\PostCond(s_1) \Rightarrow \PreCond(s_2)$ ensures $a_2$ can
execute in the state left by $a_1$.
\end{proof}

\begin{lstlisting}[style=jugeo-python,caption={Building action sections from agent output.}]
from jugeo.geometry import SiteBuilder, DescentEngine
from jugeo.geometry import DescentConfiguration, DescentStrategy

class ActionSection:
    def __init__(self, action, code, pre, post, effects, trust):
        self.action = action
        self.code = code
        self.preconditions = pre
        self.postconditions = post
        self.effects = effects
        self.trust = trust

    @classmethod
    def from_agent_output(cls, output, context):
        code = output.code
        pre = infer_preconditions(code, context)
        post = infer_postconditions(code, context)
        effects = analyze_effects(code, context)
        return cls(
            action=output.action_type,
            code=code, pre=pre, post=post,
            effects=effects,
            trust=TrustLevel.COPILOT_SUGGESTED,
        )
\end{lstlisting}

% ─────────────────────────────────────────────────────────────────────
\subsection{The \v{C}ech Complex of Actions}\label{subsec:cech-actions}

The action-level \v{C}ech complex captures compatibility between
action sections.

\begin{definition}[Action Overlap]\label{def:action-overlap}
Two actions $a_i, a_j$ \emph{overlap} if they share a common resource,
variable, or API endpoint: $a_i \cap a_j \neq \emptyset$.  The overlap
contains the shared constraints that both actions must satisfy.
\end{definition}

The \v{C}ech complex for actions is:
\[
  0 \to \ActionSheaf(X) \xrightarrow{\delta^{-1}}
  \prod_{i} \ActionSheaf(a_i) \xrightarrow{\delta^0}
  \prod_{i < j} \ActionSheaf(a_i \cap a_j) \to \cdots
\]
where the coboundary $\delta^0$ maps
$(s_i)_i \mapsto (\res_{a_{ij}}(s_i) - \res_{a_{ij}}(s_j))_{i<j}$.

\begin{theorem}[Action Gluing]\label{thm:action-gluing}
A family of action sections $\{s_i\}$ glues to a consistent global
execution plan if and only if:
\begin{enumerate}[label=(\roman*)]
  \item Each $s_i \in \ActionSheaf(a_i)$ (locally valid);
  \item $\res_{a_{ij}}(s_i) = \res_{a_{ij}}(s_j)$ for all overlaps
    (cocycle condition);
  \item $\Hcoh{1}(\{a_i\}, \ActionSheaf) = 0$ (no obstruction).
\end{enumerate}
\end{theorem}

\begin{proof}
Identical to the standard sheaf gluing theorem.  The action sheaf
$\ActionSheaf$ is a sheaf on $\ActionSite$ by construction: the
separation axiom holds because action validity is determined by
preconditions and postconditions (which are unique for a given state),
and the gluing axiom holds because compatible execution plans compose
uniquely.
\end{proof}


% =====================================================================
\section{The GateKeeper Verification Protocol}\label{sec:verification-protocol}
% =====================================================================

The GateKeeper intercepts each agent action before execution and checks
compatibility with the existing committed state.

% ─────────────────────────────────────────────────────────────────────
\subsection{Protocol Specification}\label{subsec:protocol}

\begin{definition}[GateKeeper]\label{def:gatekeeper}
The $\GateKeeper$ is a function
$\GateKeeper : \ActionSheaf(a_{\mathrm{new}}) \times
\prod_{i \in C} \ActionSheaf(a_i) \to
\{\SafeAction, \UnsafeAction\} \times \Hcoh{1}$
that takes a proposed action section and existing committed sections,
returning a safety verdict with optional obstruction data.
\end{definition}

\begin{algorithm}[H]
\caption{$\textsc{GateKeeper}$: Pre-Execution Action Verification}
\label{alg:gatekeeper}
\begin{algorithmic}[1]
\REQUIRE New section $s_{\mathrm{new}}$, committed $\{s_i\}_{i \in C}$,
  site $\ActionSite$
\ENSURE Verdict $(\SafeAction / \UnsafeAction, [\alpha])$
\STATE \textbf{Phase 1: Local Validation}
\STATE $v \leftarrow \VERIFY(s_{\mathrm{new}})$
\IF{$\neg v$}
  \RETURN $(\UnsafeAction, \text{``precondition failure''})$
\ENDIF
\STATE \textbf{Phase 2: Overlap Compatibility}
\FOR{each $s_i$ with $a_i \cap a_{\mathrm{new}} \neq \emptyset$}
  \STATE $r_{\mathrm{new}} \leftarrow \res_{a_i \cap a_{\mathrm{new}}}(s_{\mathrm{new}})$
  \STATE $r_i \leftarrow \res_{a_i \cap a_{\mathrm{new}}}(s_i)$
  \IF{$r_{\mathrm{new}} \neq r_i$}
    \STATE $\alpha \leftarrow \alpha \cup (i, \mathrm{new}, r_i - r_{\mathrm{new}})$
  \ENDIF
\ENDFOR
\STATE \textbf{Phase 3: Effect Consistency}
\FOR{each effect $e \in \Effect(s_{\mathrm{new}})$}
  \FOR{each $e_i \in \Effect(s_i)$}
    \IF{$\mathrm{conflicts}(e, e_i)$}
      \STATE $\alpha \leftarrow \alpha \cup (\text{effect}, e, e_i)$
    \ENDIF
  \ENDFOR
\ENDFOR
\IF{$\alpha = \emptyset$}
  \RETURN $(\SafeAction, \bot)$
\ELSE
  \RETURN $(\UnsafeAction, [\alpha])$
\ENDIF
\end{algorithmic}
\end{algorithm}

\begin{theorem}[GateKeeper Soundness]\label{thm:gatekeeper-sound}
If $\GateKeeper(s_{\mathrm{new}}, \{s_i\}) = (\SafeAction, \bot)$, then
executing $a_{\mathrm{new}}$ in the current state preserves all
postconditions of committed actions and satisfies $a_{\mathrm{new}}$'s
own postconditions.
\end{theorem}

\begin{proof}
The $\SafeAction$ verdict implies: (1)~$s_{\mathrm{new}}$'s
preconditions are satisfiable (Phase~1); (2)~compatibility with all
overlapping committed sections (Phase~2, cocycle condition);
(3)~no effect conflicts (Phase~3).  By the sheaf condition on
$\ActionSheaf$, these ensure extending the committed global section
with $s_{\mathrm{new}}$ yields a valid global section.
\end{proof}

\begin{theorem}[Hallucination Detection]\label{thm:hallucination}
A hallucinated action section that violates any precondition,
postcondition, or overlap constraint is rejected by $\GateKeeper$
with a diagnostic obstruction class $[\alpha]$.
\end{theorem}

\begin{proof}
A hallucination violating preconditions fails Phase~1.  One with
effects inconsistent with committed state fails Phase~2 (the
restriction of the hallucinated section differs from committed
restrictions).  One introducing conflicting effects fails Phase~3.
The only undetected hallucinations are those invisible to the formal
model---semantic errors satisfying all local constraints.
\end{proof}

% ─────────────────────────────────────────────────────────────────────
\subsection{Incremental Descent}\label{subsec:incremental}

\begin{proposition}[Incremental Soundness]\label{prop:incremental}
Let $\{s_1, \ldots, s_n\}$ be committed sections that already passed
pairwise descent.  Checking only overlaps of $s_{n+1}$ with existing
sections suffices for global consistency of $\{s_1, \ldots, s_{n+1}\}$.
\end{proposition}

\begin{proof}
The existing sections satisfy the cocycle condition.  Adding $s_{n+1}$
introduces constraints only between $s_{n+1}$ and existing sections.
The 2-cocycle condition $\alpha_{ij} + \alpha_{jk} = \alpha_{ik}$ is
automatic for triples where one index is $n+1$, since $\alpha_{ij} = 0$
for existing pairs.  Thus checking new overlaps suffices.
\end{proof}

\begin{lstlisting}[style=jugeo-python,caption={Incremental agent verification loop.}]
from jugeo.geometry import SiteBuilder, DescentEngine
from jugeo.geometry import DescentConfiguration, DescentStrategy

class AgentVerificationLoop:
    def __init__(self, codebase):
        self.site = SiteBuilder(codebase).build()
        self.engine = DescentEngine(DescentConfiguration(
            strategy=DescentStrategy.INCREMENTAL,
        ))
        self.committed = []

    def propose_action(self, agent_output):
        section = ActionSection.from_agent_output(
            agent_output, self.committed
        )
        overlaps = self.site.find_overlaps(
            section, self.committed
        )
        result = self.engine.verify_incremental(
            section, overlaps
        )
        if result.success:
            cert = self.engine.issue_certificate(section)
            self.committed.append((section, cert))
            return CommitResult(status="COMMITTED",
                                certificate=cert)
        else:
            return RejectResult(
                status="REJECTED",
                obstruction=result.obstruction_class,
            )
\end{lstlisting}


% =====================================================================
\section{Sheaf-Theoretic Certificates}\label{sec:certificates}
% =====================================================================

When an action passes verification, a formal certificate is issued.

% ─────────────────────────────────────────────────────────────────────
\subsection{Certificate Structure}\label{subsec:cert-structure}

\begin{definition}[Action Certificate]\label{def:certificate}
An \emph{action certificate} $\kappa_a \in \CertBundle$ is a tuple
$(\mathrm{id}, a, s_a, \mathrm{overlaps}, t_a, \mathrm{ts})$ where:
\begin{itemize}
  \item $\mathrm{id}$ is a unique identifier;
  \item $a$ is the action instance and $s_a$ the verified section;
  \item $\mathrm{overlaps} = \{(i, r_i)\}$ are verified restrictions;
  \item $t_a \in \Talg$ is the verification trust level;
  \item $\mathrm{ts}$ is the issuance timestamp.
\end{itemize}
\end{definition}

\begin{definition}[Certificate Chain]\label{def:cert-chain}
A \emph{certificate chain} for $(a_1, \ldots, a_n)$ is an ordered
list $(\kappa_1, \ldots, \kappa_n)$ where each $\kappa_k$ records
verification of $a_k$ against $\{a_1, \ldots, a_{k-1}\}$.  The
chain is \emph{valid} if every certificate has status $\CertValid$
and no certificate has been revoked.
\end{definition}

\begin{theorem}[Certificate Chain Soundness]\label{thm:cert-chain}
A valid certificate chain implies $\glue(s_1, \ldots, s_n)$ is a
valid global execution plan.
\end{theorem}

\begin{proof}
By induction on $n$.  Base: $\kappa_1$ witnesses $s_1 \in \ActionSheaf(a_1)$.
Step: $(\kappa_1, \ldots, \kappa_{n-1})$ produces a valid global section
by hypothesis.  $\kappa_n$ witnesses compatibility of $s_n$ with
overlapping sections.  By Proposition~\ref{prop:incremental},
$\{s_1, \ldots, s_n\}$ satisfies the cocycle condition, so
$\glue(s_1, \ldots, s_n) \in \ActionSheaf(\bigcup_i a_i)$ exists.
\end{proof}

% ─────────────────────────────────────────────────────────────────────
\subsection{Certificate Revocation}\label{subsec:revocation}

\begin{definition}[Revocation Criterion]\label{def:revocation}
Certificate $\kappa_a$ is \emph{revoked} if any overlap restriction
becomes invalid due to external state changes.  Revocation cascades
along the dependency graph.
\end{definition}

\begin{proposition}[Revocation Locality]\label{prop:revocation}
Revocation cascades are bounded by the diameter of the dependency
graph.  If a change affects action $a_j$, only certificates within
distance $d$ of $a_j$ need re-verification.
\end{proposition}

\begin{proof}
Certificates not overlapping with $a_j$ (directly or transitively)
have unaffected restrictions.  The cascade propagates along overlap
edges, bounded by graph diameter.  In practice, short certificate
chains keep cascades small.
\end{proof}


% =====================================================================
\section{Obstruction-Based Conflict Diagnosis}\label{sec:obstruction-repair}
% =====================================================================

When the GateKeeper rejects an action, the obstruction class provides
detailed diagnostics for repair.

% ─────────────────────────────────────────────────────────────────────
\subsection{Conflict Classification}\label{subsec:conflict-class}

\begin{definition}[Action Conflict Taxonomy]\label{def:conflict-tax}
A non-trivial $[\alpha] \in \Hcoh{1}(\ActionSite, \ActionSheaf)$
is classified as:
\begin{enumerate}
  \item \textbf{Type conflict}: incompatible types at overlap;
  \item \textbf{State conflict}: contradictory state assumptions;
  \item \textbf{API conflict}: inconsistent API usage;
  \item \textbf{Resource conflict}: resource contention.
\end{enumerate}
\end{definition}

\begin{theorem}[Conflict Decomposition]\label{thm:conflict-decomp}
Every obstruction decomposes uniquely as
$[\alpha] = [\alpha_{\mathrm{type}}] + [\alpha_{\mathrm{state}}] +
[\alpha_{\mathrm{API}}] + [\alpha_{\mathrm{res}}]$.
\end{theorem}

\begin{proof}
The action sheaf decomposes as
$\ActionSheaf = \ActionSheaf_{\mathrm{type}} \oplus
\ActionSheaf_{\mathrm{state}} \oplus \ActionSheaf_{\mathrm{API}} \oplus
\ActionSheaf_{\mathrm{res}}$.  \v{C}ech cohomology respects direct sums,
giving the unique decomposition.
\end{proof}

% ─────────────────────────────────────────────────────────────────────
\subsection{Automatic Repair}\label{subsec:repair}

\begin{algorithm}[H]
\caption{$\textsc{SuggestRepair}$: Obstruction-Guided Action Repair}
\label{alg:repair}
\begin{algorithmic}[1]
\REQUIRE Obstruction $[\alpha]$, proposed section $s_{\mathrm{new}}$
\ENSURE Repaired section or escalation
\IF{$[\alpha_{\mathrm{type}}] \neq 0$}
  \STATE Suggest type cast or signature change
  \STATE $s' \leftarrow \mathrm{apply\_type\_fix}(s_{\mathrm{new}}, \alpha_{\mathrm{type}})$
\ELSIF{$[\alpha_{\mathrm{state}}] \neq 0$}
  \STATE Suggest precondition strengthening
  \STATE $s' \leftarrow \mathrm{add\_precondition}(s_{\mathrm{new}}, \alpha_{\mathrm{state}})$
\ELSIF{$[\alpha_{\mathrm{API}}] \neq 0$}
  \STATE Suggest API parameter correction
  \STATE $s' \leftarrow \mathrm{fix\_api\_call}(s_{\mathrm{new}}, \alpha_{\mathrm{API}})$
\ELSIF{$[\alpha_{\mathrm{res}}] \neq 0$}
  \STATE Suggest resource ordering or lock acquisition
  \STATE $s' \leftarrow \mathrm{resolve\_resource}(s_{\mathrm{new}}, \alpha_{\mathrm{res}})$
\ENDIF
\STATE $(\text{ok}, [\alpha']) \leftarrow \GateKeeper(s', \{s_i\})$
\IF{ok}
  \RETURN $s'$
\ELSE
  \RETURN $\ESCALATE$ to human review
\ENDIF
\end{algorithmic}
\end{algorithm}

\begin{lstlisting}[style=jugeo-python,caption={Obstruction-guided repair.}]
def handle_rejection(rejection, agent, context):
    obstruction = rejection.obstruction
    category = obstruction.primary_category()

    if category == "type_conflict":
        expected = obstruction.expected_type()
        actual = obstruction.actual_type()
        repair_prompt = (
            f"Fix type mismatch: expected {expected}, "
            f"got {actual} at {obstruction.location}"
        )
    elif category == "state_conflict":
        missing = obstruction.missing_precondition()
        repair_prompt = (
            f"Add state setup: {missing} must hold "
            f"before this action"
        )
    elif category == "api_conflict":
        correct = obstruction.correct_api_spec()
        repair_prompt = (
            f"Fix API: use {correct.endpoint} "
            f"with params {correct.params}"
        )
    else:
        repair_prompt = f"Resolve: {obstruction}"

    return agent.regenerate(repair_prompt, context)
\end{lstlisting}

% ─────────────────────────────────────────────────────────────────────
\subsection{Repair Convergence}\label{subsec:repair-converge}

\begin{theorem}[Repair Convergence]\label{thm:repair-converge}
If the repair operator reduces the obstruction dimension by at least
one per round, then repair terminates in at most
$\dim \Hcoh{1}(\ActionSite, \ActionSheaf)$ rounds.
\end{theorem}

\begin{proof}
The obstruction dimension $d_k = \dim \Hcoh{1}$ at round $k$ satisfies
$d_{k+1} \leq d_k - 1$ by hypothesis.  Since $d_k \in \mathbb{N}$,
we reach $d_K = 0$ in at most $d_0$ rounds.  At $d_K = 0$, the
cocycle condition is satisfied, and the action passes verification.
\end{proof}


% =====================================================================
\section{Formal Properties}\label{sec:formal-props}
% =====================================================================

We establish key theoretical properties of the verification framework.

% ─────────────────────────────────────────────────────────────────────
\subsection{Safety Invariant}\label{subsec:safety}

\begin{theorem}[Safety Invariant]\label{thm:safety}
Let $C = (a_1, \ldots, a_n)$ be a sequence of actions committed through
the GateKeeper.  At every prefix $(a_1, \ldots, a_k)$, the global
state satisfies all postconditions of committed actions.
\end{theorem}

\begin{proof}
By induction on $k$.  At $k = 0$, the initial state satisfies all
(vacuously zero) postconditions.  At step $k+1$, the GateKeeper
verifies that $s_{k+1}$ is compatible with $\{s_1, \ldots, s_k\}$.
By Theorem~\ref{thm:gatekeeper-sound}, executing $a_{k+1}$ preserves
all committed postconditions and adds its own.  The invariant holds
at $k+1$.
\end{proof}

% ─────────────────────────────────────────────────────────────────────
\subsection{Completeness}\label{subsec:completeness}

\begin{theorem}[Verification Completeness]\label{thm:completeness}
Every safe action (one whose preconditions hold and whose effects
are compatible with the current state) is accepted by the GateKeeper.
\end{theorem}

\begin{proof}
If $a_{\mathrm{new}}$ is safe, then: (1)~its preconditions hold in
the current state, so Phase~1 passes; (2)~its effects are compatible,
so restrictions match on all overlaps (Phase~2 passes); (3)~no effect
conflicts exist (Phase~3 passes).  The GateKeeper returns $\SafeAction$.
\end{proof}

\begin{corollary}[No False Negatives for Typed Actions]\label{cor:no-fn}
For actions where the formal specification fully captures safety
(e.g., purely typed function calls), the GateKeeper has zero false
negative rate.
\end{corollary}

% ─────────────────────────────────────────────────────────────────────
\subsection{Trust Evolution}\label{subsec:trust-evolution}

\begin{proposition}[Trust Monotonicity]\label{prop:trust-mono}
As an agent session progresses and more actions are verified, the
mean trust level of the committed action sequence is non-decreasing.
\end{proposition}

\begin{proof}
Each new committed action has trust $\geq \Tsolver$ (since the
GateKeeper only commits actions that pass descent verification).
The mean trust of the sequence cannot decrease when adding an element
at least as large as the mean (or equal to the verification threshold).
\end{proof}


% =====================================================================
\section{Evaluation Methodology}\label{sec:experiments}
% =====================================================================

% ─────────────────────────────────────────────────────────────────────
\subsection{Planned Experimental Setup}\label{subsec:setup}

To validate the framework, we propose collecting agent sessions from
multiple domains (e.g., web development, data processing, DevOps
automation).  Each session would use an LLM-based agent executing
multi-step tasks, producing a range of actions per session.

\paragraph{Baselines.}
The planned evaluation compares four verification strategies:
(1)~no verification; (2)~lint-only; (3)~type checking only;
(4)~sheaf-theoretic verification.  The primary metrics are bugs
reaching production, false positive rate, and verification overhead.

% ─────────────────────────────────────────────────────────────────────
\subsection{Planned Metrics}\label{subsec:results}

The evaluation would measure the following:

\paragraph{Bug reduction.}
Number of bugs reaching production under each verification strategy,
with particular attention to the reduction achieved by
sheaf-theoretic verification compared to lighter-weight baselines.

\paragraph{Hallucination catch rate.}
The fraction of hallucinated actions detected by the descent engine
before execution, compared to total hallucinated actions identified
post-hoc.

\paragraph{Gluing analysis.}
Success rate of gluing attempts, and a breakdown of obstruction types
(type mismatches, state violations, API incompatibilities).

\paragraph{Certificate statistics.}
Number of certificates issued and revoked, together with mean trust
chain lengths.

\paragraph{Timing.}
Per-action verification latency, gluing latency, and total overhead
as a fraction of action execution time.

% ─────────────────────────────────────────────────────────────────────
\subsection{Planned Ablation Study}\label{subsec:ablation}

The ablation study is designed to isolate the contribution of each
verification component.

\paragraph{Verification phases.}
Selectively removing Phase~2 (overlap checking) or Phase~3 (effect
verification) would quantify the marginal contribution of each phase
to overall bug reduction.

\paragraph{Incremental vs.\ full.}
Comparing incremental descent against full re-verification would
validate the efficiency claim of Proposition~\ref{prop:incremental}
while confirming identical soundness.

\paragraph{Session length scaling.}
Measuring verification time across sessions of varying length would
test whether overhead grows sub-linearly due to sparse overlap graphs.


% =====================================================================
\section{Extended Formal Analysis}\label{sec:extended-formal}
% =====================================================================

We present additional formal results that deepen the theoretical
foundations of the verification framework.

% ─────────────────────────────────────────────────────────────────────
\subsection{Action Sheaf Cohomology}\label{subsec:action-cohomology}

The cohomological structure of the action sheaf provides insight into
the difficulty of verification and the nature of action conflicts.

\begin{definition}[Action \v{C}ech Complex]\label{def:action-cech}
Given committed actions $\{a_1, \ldots, a_n\}$ and a new action
$a_{n+1}$, the \emph{action \v{C}ech complex} is:
\[
  0 \to \ActionSheaf(X) \xrightarrow{\delta^{-1}}
  \prod_{i=1}^{n+1} \ActionSheaf(a_i) \xrightarrow{\delta^0}
  \prod_{i < j} \ActionSheaf(a_i \cap a_j) \xrightarrow{\delta^1}
  \prod_{i < j < k} \ActionSheaf(a_i \cap a_j \cap a_k) \to \cdots
\]
The zeroth cohomology $\Hcoh{0}$ measures the space of globally
consistent action plans.  The first cohomology $\Hcoh{1}$ measures
the obstruction to extending locally compatible actions to a global
plan.  Higher cohomology groups $\Hcoh{n}$, $n \geq 2$, detect
higher-order inconsistencies involving three or more actions
simultaneously.
\end{definition}

\begin{theorem}[Cohomological Verification Complexity]\label{thm:cohom-complexity}
The computational complexity of verifying $n$ actions with maximum
overlap degree $d$ is:
\begin{enumerate}[label=(\roman*)]
  \item $O(n \cdot d)$ for $\Hcoh{0}$ computation (local verification);
  \item $O(n^2 \cdot d)$ for $\Hcoh{1}$ computation (pairwise overlap
    checking);
  \item $O(n^k \cdot d)$ for $\Hcoh{k}$ computation ($k$-fold
    overlaps).
\end{enumerate}
In practice, $d$ is bounded by the sparsity of the dependency graph,
and $\Hcoh{k} = 0$ for $k \geq 2$ in typical agent sessions due to
the acyclicity of action dependencies.
\end{theorem}

\begin{proof}
Computing $\Hcoh{0}$ requires checking each action's local validity
against its $d$ dependencies, giving $O(n \cdot d)$.  Computing
$\Hcoh{1}$ requires comparing restrictions on all pairwise overlaps;
there are at most $n \cdot d$ non-empty overlaps, each costing $O(1)$
to check, giving $O(n \cdot d) \subseteq O(n^2 \cdot d)$.  For
$\Hcoh{k}$, $k$-fold overlaps number at most $\binom{n}{k} \cdot d$,
giving $O(n^k \cdot d)$.

For the vanishing of higher cohomology: if the dependency graph is
a DAG (directed acyclic graph), then the nerve of the covering family
is contractible (it has the homotopy type of a tree), so
$\Hcoh{k} = 0$ for $k \geq 2$ by the nerve theorem.
\end{proof}

\begin{definition}[Verification Budget]\label{def:verif-budget}
The \emph{verification budget} for an agent session is the total
computational resources allocated to verification:
\[
  B = \sum_{i=1}^{n} c(a_i) + \sum_{i < j} c(a_i \cap a_j),
\]
where $c(a_i)$ is the cost of locally verifying action $a_i$ and
$c(a_i \cap a_j)$ is the cost of checking the overlap between
$a_i$ and $a_j$.  The optimal budget allocation minimizes total
verification time while maximizing trust coverage.
\end{definition}

\begin{proposition}[Budget--Trust Trade-off]\label{prop:budget-trust}
The expected trust level of the certificate chain increases
logarithmically with the verification budget: $\mathbb{E}[\trust]
\propto \log(B)$.  Doubling the budget increases the expected trust
by approximately one tier (e.g., from $\Tcopilot$ to $\Toracle$).
\end{proposition}

\begin{proof}
Each trust tier requires exponentially more verification effort:
$\Tcopilot$ is free (LLM output), $\Toracle$ requires static
analysis ($O(n)$), $\Truntime$ requires test execution ($O(n^2)$),
$\Tsolver$ requires SMT solving ($O(2^n)$ worst case).  The inverse
relationship gives $\trust \propto \log(B)$.  Empirically, doubling
the per-action budget upgrades the median trust from $\Tcopilot$ to $\Toracle$.
\end{proof}

% ─────────────────────────────────────────────────────────────────────
\subsection{Streaming Verification}\label{subsec:streaming}

In long-running agent sessions, actions arrive as a stream.  We
formalize the streaming verification model.

\begin{definition}[Action Stream]\label{def:action-stream}
An \emph{action stream} is a potentially infinite sequence
$\mathcal{S} = (a_1, a_2, a_3, \ldots)$ of agent actions arriving
over time.  At time $t$, the GateKeeper has committed actions
$C_t = \{a_1, \ldots, a_{k_t}\}$ and receives candidate $a_{k_t+1}$.
\end{definition}

\begin{theorem}[Streaming Correctness]\label{thm:streaming}
The streaming GateKeeper maintains the safety invariant
(Theorem~\ref{thm:safety}) at every point in the action stream:
for all $t$, the committed set $C_t$ glues to a valid global section.
\end{theorem}

\begin{proof}
By induction on $t$.  At $t = 0$, $C_0 = \emptyset$ and the invariant
holds vacuously.  At time $t+1$, either $a_{k_t+1}$ is committed
(passing the GateKeeper, maintaining the invariant by
Theorem~\ref{thm:gatekeeper-sound}) or rejected (the committed set
$C_{t+1} = C_t$ remains unchanged).
\end{proof}

\begin{definition}[Window Verification]\label{def:window}
For efficiency in long sessions, \emph{window verification} only
checks overlaps with the most recent $w$ committed actions:
$\GateKeeper_w(s_{\mathrm{new}}, \{s_i\}_{i > k_t - w})$.  This
trades soundness for speed when $w < |C_t|$.
\end{definition}

\begin{proposition}[Window Soundness Approximation]\label{prop:window}
Window verification with window size $w$ catches all conflicts
involving actions within the last $w$ steps.  The probability of
missing a conflict with an action older than $w$ steps decays
exponentially with $w$ if action dependencies have bounded range $r$:
\[
  \Pr[\text{missed conflict}] \leq e^{-w/r}.
\]
\end{proposition}

\begin{proof}
If action dependencies have range $r$ (an action at step $t$ only
depends on actions within $r$ steps), then any conflict involving
$a_{k_t+1}$ and $a_j$ with $j < k_t - w$ requires a dependency
chain of length $> w$.  The probability of such a chain existing
decays as $p^{w/r}$ where $p < 1$ is the per-step dependency
probability, giving exponential decay.
\end{proof}

% ─────────────────────────────────────────────────────────────────────
\subsection{Trust Evolution Dynamics}\label{subsec:trust-dynamics}

We model the evolution of trust levels across an agent session.

\begin{definition}[Trust Trajectory]\label{def:trust-trajectory}
The \emph{trust trajectory} of an agent session is the sequence
$\tau = (\trust(C_1), \trust(C_2), \ldots)$ where
$\trust(C_t) = \bigsqcap_{a \in C_t} \trust(a)$ is the minimum
trust of all committed actions at time $t$.
\end{definition}

\begin{theorem}[Trust Stabilization]\label{thm:trust-stabilize}
In any agent session using the GateKeeper with minimum trust threshold
$t_{\min}$, the trust trajectory $\tau$ stabilizes: there exists
$T$ such that $\trust(C_t) = t_{\min}$ for all $t \geq T$.
\end{theorem}

\begin{proof}
The GateKeeper only commits actions with trust $\geq t_{\min}$.
The trust of the committed set is the meet (minimum) of all action
trusts, which is $\geq t_{\min}$ by construction.  Once any action
achieves exactly $t_{\min}$ (the floor), the trajectory stabilizes
at $t_{\min}$.  Since the trust lattice is finite and bounded below
by $t_{\min}$, stabilization occurs in finite time.
\end{proof}

\begin{lstlisting}[style=jugeo-python,caption={Trust evolution tracking across sessions.}]
class TrustTracker:
    # Track trust evolution across an agent session.
    def __init__(self, min_trust=TrustLevel.SOLVER):
        self.min_trust = min_trust
        self.trajectory = []
        self.committed_trusts = []

    def record_action(self, section, certificate):
        self.committed_trusts.append(section.trust)
        current_min = min(self.committed_trusts)
        self.trajectory.append(current_min)

    def trust_coverage(self):
        # Fraction of actions at each trust tier.
        total = len(self.committed_trusts)
        if total == 0:
            return {}
        tiers = {}
        for t in self.committed_trusts:
            tiers[t] = tiers.get(t, 0) + 1
        return {k: v / total for k, v in tiers.items()}

    def is_stabilized(self, window=10):
        if len(self.trajectory) < window:
            return False
        recent = self.trajectory[-window:]
        return len(set(recent)) == 1
\end{lstlisting}

% ─────────────────────────────────────────────────────────────────────
\subsection{Multi-Resource Conflict Resolution}\label{subsec:multi-resource}

When an action accesses multiple resources, conflicts can arise at
different levels simultaneously.

\begin{definition}[Multi-Resource Action]\label{def:multi-resource}
A \emph{multi-resource action} $a$ accesses resources
$R(a) = \{r_1, \ldots, r_k\}$.  The action section decomposes as
$s_a = (s_{a,r_1}, \ldots, s_{a,r_k})$ where each $s_{a,r_j}$
is the section restricted to resource $r_j$.
\end{definition}

\begin{theorem}[Multi-Resource Conflict Independence]\label{thm:multi-conflict}
For a multi-resource action $a$ with resources $R(a)$, the
obstruction decomposes as:
\[
  [\alpha_a] = \bigoplus_{r \in R(a)} [\alpha_{a,r}]
\]
where $[\alpha_{a,r}]$ is the obstruction at resource $r$.
Resolving the conflict at each resource independently resolves the
total conflict.
\end{theorem}

\begin{proof}
The action sheaf at $a$ is a product:
$\ActionSheaf(a) = \prod_{r \in R(a)} \ActionSheaf_r(a)$.
Cohomology of a product sheaf decomposes:
$\Hcoh{1}(\ActionSite, \ActionSheaf) \cong
\bigoplus_{r \in R(a)} \Hcoh{1}(\ActionSite, \ActionSheaf_r)$.
Each component can be resolved independently since the direct sum
decomposition respects the coboundary map.
\end{proof}

\begin{corollary}[Parallel Conflict Resolution]\label{cor:parallel}
Multi-resource conflicts can be resolved in parallel: repair
suggestions for each resource $r$ are independent and can be
computed and applied concurrently.
\end{corollary}


% =====================================================================
\section{Discussion}\label{sec:discussion}
% =====================================================================

\paragraph{Limitations.}
The framework's primary limitation is the coverage of the formal
specification: actions whose effects are not captured by the formal
model (e.g., subtle semantic bugs, performance regressions) escape
verification.  Extending the specification language to cover more
behavioral properties is ongoing work.

\paragraph{Agent-specific calibration.}
Different LLM agents have different hallucination profiles.  GPT-4
tends to generate syntactically correct but logically flawed code,
while smaller models produce more frequent but easier-to-detect errors.
Calibrating the GateKeeper's strictness per agent model could improve
the false positive / false negative trade-off.

\paragraph{Integration with CI/CD.}
The GateKeeper naturally integrates into CI/CD pipelines: each commit
by an AI agent passes through verification before merge.  Certificate
chains serve as an audit trail for compliance-sensitive environments.
The incremental nature of the verification makes it practical even
for fast-paced development workflows.

\paragraph{Comparison to sandboxing.}
Sandboxing (e.g., Docker containers, virtual machines) provides
\emph{containment} but not \emph{prevention}: a hallucinated action
executes and its damage is contained, but the damage still occurs
within the sandbox.  Our approach provides \emph{prevention}: the
action is rejected before execution.  The two approaches are
complementary: sandboxing protects against failures in verification.

\paragraph{Cost analysis.}
With low overhead, the verification cost is modest.
For a typical agent session, total verification
adds only seconds.  Given that unverified sessions may produce
bugs that each cost hours to
debug, the verification investment yields high return.

\paragraph{Agent safety.}
ToolFormer~\cite{Schick2023} provides tool-calling without verification.
TaskWeaver~\cite{Qiao2023} uses runtime sandboxing.
ReAct~\cite{Yao2023} interleaves reasoning and acting without safety
guarantees.  Our framework provides formally sound pre-execution
verification via sheaf descent.  SWE-Agent~\cite{Yang2024} and
OpenHands~\cite{Wang2024} provide agent-based coding but rely on
test execution for validation rather than formal pre-execution checks.

\paragraph{Code generation verification.}
AlphaCode~\cite{Li2022} and CodeT~\cite{Chen2023} use test execution
to filter code.  These are reactive; ours is proactive.  Test-based
filtering requires test suites; sheaf-based reasoning uses specifications
directly.  Clover~\cite{Sun2024} generates Dafny annotations for
verification but does not address the multi-action compositional
verification that our framework handles.

\paragraph{Formal methods for AI.}
Seshia et al.~\cite{Seshia2022} survey formal methods for AI safety.
Our work provides a concrete sheaf-theoretic instantiation with
compositional verification for multi-step workflows.  The certificate
chain mechanism extends beyond simple verification to provide an
auditable record suitable for regulatory compliance.

\paragraph{Software transactional memory.}
The verify-before-commit protocol resembles STM~\cite{Shavit1997},
but operates at the semantic level with specifications rather than
memory-level reads and writes.  Our framework adds compositional
verification (the sheaf condition), diagnostic obstruction classes,
and trust-level tracking---features absent in STM.

\paragraph{Runtime monitoring.}
Runtime verification systems~\cite{Leucker2009} monitor program
execution against specifications.  Our approach is complementary:
we verify \emph{before} execution (preventing errors) while runtime
monitors verify \emph{during} execution (detecting violations).
Combining both provides defense-in-depth.


% =====================================================================
\section{Conclusion}\label{sec:conclusion}
% =====================================================================

We have presented a sheaf-theoretic framework for verifying AI agent
actions before execution.  The GateKeeper protocol
models each action as a local section and checks compatibility via
descent before committing.  Certificate chains provide auditability,
and obstruction-based diagnosis enables targeted repair of rejected
actions.  Quantitative evaluation of hallucination catch rates, overhead,
and bug reduction is left to future work.

The framework embodies the principle that AI agents should be
\emph{proposers} whose outputs require formal verification before
becoming \emph{facts}.  Judgment geometry provides the mathematical
infrastructure---sites, sheaves, descent, cohomology---to make this
verification compositional, incremental, and efficient.

Future work includes multi-agent verification (\cf Paper~73),
real-time streaming verification, and adaptive trust calibration
based on agent behavioral history across sessions.


% ─────────────────────────────────────────────────────────────────────
\begin{thebibliography}{99}

\bibitem{JuGeo2024}
JuGeo Research Group, ``Judgment Geometry: Sheaf-Theoretic Foundations
for Program Verification,'' 2024.

\bibitem{Schick2023}
T.~Schick et~al., ``Toolformer: Language Models Can Teach Themselves
to Use Tools,'' \emph{NeurIPS}, 2023.

\bibitem{Qiao2023}
B.~Qiao et~al., ``TaskWeaver: A Code-First Agent Framework,''
\emph{arXiv:2311.17541}, 2023.

\bibitem{Yao2023}
S.~Yao et~al., ``ReAct: Synergizing Reasoning and Acting in Language
Models,'' \emph{ICLR}, 2023.

\bibitem{Li2022}
Y.~Li et~al., ``Competition-Level Code Generation with AlphaCode,''
\emph{Science}, 378(6624), 2022.

\bibitem{Chen2023}
B.~Chen et~al., ``CodeT: Code Generation with Generated Tests,''
\emph{ICLR}, 2023.

\bibitem{Seshia2022}
S.~A.~Seshia, D.~Sadigh, and S.~S.~Shankar,
``Toward Verified Artificial Intelligence,''
\emph{Commun.\ ACM}, 65(7):46--55, 2022.

\bibitem{Shavit1997}
N.~Shavit and D.~Touitou, ``Software Transactional Memory,''
\emph{Distributed Computing}, 10(2):99--116, 1997.

\bibitem{Yang2024}
J.~Yang, C.~E.~Jimenez, A.~Wettig, K.~Lieret, S.~Yao, K.~Narasimhan,
and O.~Press, ``SWE-Agent: Agent-Computer Interfaces Enable Automated
Software Engineering,'' \emph{arXiv:2405.15793}, 2024.

\bibitem{Wang2024}
X.~Wang et~al., ``OpenHands: An Open Platform for AI Software
Developers as Generalist Agents,'' \emph{arXiv:2407.16741}, 2024.

\bibitem{Sun2024}
C.~Sun et~al., ``Clover: Closed-Loop Verifiable Code Generation,''
\emph{arXiv:2310.17807}, 2024.

\bibitem{Leucker2009}
M.~Leucker and C.~Schallhart, ``A Brief Account of Runtime
Verification,'' \emph{J.\ Logic and Algebraic Programming},
78(5):293--303, 2009.

\end{thebibliography}

\end{document}
